\section*{Chapter \ref{ch:g12:ideas-institutions-growth} --- Ideas, Institutions and Growth}
% figdata: figdata/grade-12/06-ideas-institutions-growth/out/ (island, intensity, intensity_key, patent_key); the medicine is invented.

\begin{solution}{exo:g12:ideas-institutions-growth:1}
Recipe: nonrival. Cake: rival. Program: nonrival (copies cost almost nothing). Truck: rival.
\end{solution}

\begin{solution}{exo:g12:ideas-institutions-growth:2}
Theorem: nonrival. Train seat: rival. Song: nonrival. Vaccine dose: rival. The formula: nonrival.
\end{solution}

\begin{solution}{exo:g12:ideas-institutions-growth:3}
1950: Haiti 1790, Dominican Republic 1637 dollars. 2022: Haiti 1546, Dominican Republic 17\,714. Ratio $17\,714/1546 \approx 11.5$.
\end{solution}

\begin{solution}{exo:g12:ideas-institutions-growth:4}
In the 1950s: by 1960 the Dominican Republic was ahead (2075 against 1809). In 1980: 3781 against 2063.
\end{solution}

\begin{solution}{exo:g12:ideas-institutions-growth:5}
\omterm{def:g11:marginal-reasoning:marginal-benefit}{Marginal revenue} $50 - 2Q = 10$ gives $Q = 20$, $P = 30$; profit $(30 - 10) \times 20 = 400$; \omterm{def:g10:what-prices-do:dwl}{deadweight loss} $\tfrac12 \times 20 \times 20 = 200$ (competitive quantity 40).
\end{solution}

\begin{solution}{exo:g12:ideas-institutions-growth:6}
$100 \times (1 - 1.04^{-20})/0.04 \approx 100 \times 13.590 = 1359$.
\end{solution}

\begin{solution}{exo:g12:ideas-institutions-growth:7}
$80 \times 12.462 \approx 997$: the prize of 1000 is worth slightly more.
\end{solution}

\begin{solution}{exo:g12:ideas-institutions-growth:8}
$1600 \times (14.094 - 12.462) \approx 2611$ thousand dalens. Society gives up five more years of the \omterm{def:g10:what-prices-do:dwl}{deadweight loss}, worth $800 \times 1.632 \approx 1305$, and buyers pay the
monopoly price for five years longer.
\end{solution}

\begin{solution}{exo:g12:ideas-institutions-growth:9}
$(4.24/6.72)^{1/21} - 1 \approx -2.2\,\%$ a year.
\end{solution}

\begin{solution}{exo:g12:ideas-institutions-growth:10}
10.91 in 2000 and 6.30 in 2021: a fall of about 42\,\%, against 25\,\% for the world (6.05 to 4.54).
\end{solution}

\begin{solution}{exo:g12:ideas-institutions-growth:11}
$1/(1 + e^{-x}) = p$ gives $x = \ln(p/(1 - p))$. 10\,\%: $0.3(t - 20) = -\ln 9$, $t \approx 12.7$; 50\,\%: $t = 20$; 90\,\%: $t \approx 27.3$. The second country reaches 50\,\% at
$t = 45$.
\end{solution}

\begin{solution}{exo:g12:ideas-institutions-growth:12}
The toy model predicts that rising research effort raises growth. Research effort has risen substantially for decades without faster productivity growth, and keeping up
the doubling of chip density needs more than 18 times the researchers of the early 1970s. The assumption that fails is a constant $\theta$: each researcher finds fewer
new \omterm{def:g12:ideas-institutions-growth:idea}{ideas} over time.
\end{solution}

\begin{solution}{exo:g12:ideas-institutions-growth:13}
No: it shows two countries with a similar start and different paths. Other explanations: different histories, political
instability, natural disasters, differences in schooling or in trade. To separate them one needs many cases, or a natural experiment in which the rules differ while the
other factors do not, as with the countries split by history in the grades 6--9 volume.
\end{solution}

\begin{solution}{exo:g12:ideas-institutions-growth:14}
$R = 10$: 1\,\% a year; $R = 20$: 2\,\%. If $\theta$ halves while $R$ doubles, growth stays at 1\,\%: research effort rises with no faster growth, as in the finding that \omterm{def:g12:ideas-institutions-growth:idea}{ideas}
are getting harder to find.
\end{solution}

\begin{solution}{exo:g12:ideas-institutions-growth:15}
Workers and owners of shrinking and closing firms lose, at least for a time: in one quarter, 7.3 million American jobs were lost while 7.4 million were created. The gains go
to the workers and owners of expanding firms and to buyers, through lower prices and better products. When the lost jobs are concentrated among some groups or places,
the spread of incomes can widen, which the measures of inequality would show, while average income rises.
\end{solution}

\begin{solution}{pb:g12:ideas-institutions-growth:1}
\textbf{1.} \omterm{def:g11:marginal-reasoning:marginal-benefit}{Marginal revenue} $100 - 2Q = 20$: $Q = 40$, $P = 60$.

\textbf{2.} $(60 - 20) \times 40 = 1600$ thousand dalens a year.

\textbf{3.} \omterm{def:g10:surplus-and-efficiency:consumer-surplus}{Consumer surplus} $\tfrac12 \times 40 \times 40 = 800$; \omterm{def:g10:what-prices-do:dwl}{deadweight loss} $\tfrac12 \times 40 \times 40 = 800$ (the competitive quantity is 80).

\textbf{4.} $1600 \times (1 - 1.05^{-20})/0.05 = 1600 \times 12.462 \approx 19\,940$.

\textbf{5.} Price 20, quantity 80, \omterm{def:g10:surplus-and-efficiency:consumer-surplus}{consumer surplus} $\tfrac12 \times 80 \times 80 = 3200$.

\textbf{6.} 19\,940 thousand dalens.

\textbf{7.} At the cost of 20: anyone may make it, and competition brings the price down to cost.

\textbf{8.} $800 \times 12.462 \approx 9970$.

\textbf{9.} \omterm{def:g10:surplus-and-efficiency:consumer-surplus}{Consumer surplus} rises from 800 to 3200 a year, by 2400: $2400 \times 12.462 \approx 29\,910$. Paying the prize of 19\,940 out of it leaves 9970, the \omterm{def:g10:what-prices-do:dwl}{deadweight
loss} avoided.

\textbf{10.} Yes under both: 19\,940 exceeds 15\,000.

\textbf{11.} 12\,000 is below the cost of 15\,000: no one develops the medicine, and the whole surplus is lost.

\textbf{12.} The inventor gains $30\,000 - 19\,940 = 10\,060$ more than the patent would have given, paid by taxpayers; the medicine is still sold at cost.

\textbf{13.} To set the prize, the government must estimate demand, costs and the value of the result before the medicine exists; with a patent, the market reveals the
value afterwards.

\textbf{14.} 30\,000 is spent for one medicine worth 15\,000 to develop. The loser's spending duplicates the winner's, though the race may bring the medicine sooner.

\textbf{15.} A prize needs a result that can be checked without dispute; when judging it takes years, the reward is delayed and uncertain, which weakens the incentive.

\textbf{16.} $1600 \times 14.094 \approx 22\,550$.

\textbf{17.} Only the 25-year patent: 22\,550 covers 21\,000, while 19\,940 does not.

\textbf{18.} $800 \times (14.094 - 12.462) \approx 1305$.

\textbf{19.} A longer patent raises the reward, so more costly medicines become worth developing; but each is sold at the monopoly price for longer, so buyers lose more
of the surplus of each.

\textbf{20.} A prize of about 19\,940 thousand dalens, paid now, gives the inventor the same \omterm{def:g12:present-value:present-value}{present value} as the 20-year patent, and it avoids a \omterm{def:g10:what-prices-do:dwl}{deadweight loss} worth about
9970 thousand dalens in \omterm{def:g12:present-value:present-value}{present value}.
\end{solution}
